% TOP_PROBLEM: {"id": 30006475, "title": "Torsion-Anomalous Conjecture for Semiabelian Varieties", "category_id": 5, "category": {"id": 5, "name": "algebraic_geometry", "display_name": "Algebraic Geometry", "description": "Geometric objects defined by polynomial equations.", "slug": "algebraic-geometry", "order_index": 5, "created_at": "2026-07-31T15:26:25.671Z"}, "status": "open"}
% FIRST_POSTED: 2026-10-06
% ATTEMPT_STATUS: unresolved
% OPENAI_CLAIM: {"schema_version":1,"source_repo":"https://github.com/openai/math","source_commit":"adc7f1241b42e322a6451854ab7e4b4c146bf78a","release_date":"2026-10-06","match":"partial","resolution":"partial","independently_reviewed":false,"second_pass":"n/a","families":[{"family":"016","title":"Zilber–Pink in abelian varieties and the Siegel threefold","subject":"Number theory","manuscripts":[{"title":"The abelian Zilber–Pink conjecture","pdf_path":"preprints/The-Abelian-Zilber-Pink-Conjecture-September-24-2026/paper.pdf","readme_path":"preprints/The-Abelian-Zilber-Pink-Conjecture-September-24-2026/README.md","date":"2026-09-24","theorem_ref":"thm:main (sections/01-introduction.tex:43)","statement_sources":[{"path":"preprints/The-Abelian-Zilber-Pink-Conjecture-September-24-2026/build/sections/01-introduction.tex","line":43,"label":"thm:main"}]}],"lean":null,"reasoning_trace":null}]}
\documentclass[11pt]{article}
\usepackage[margin=25mm]{geometry}
\usepackage{fontspec}
\setmainfont{Latin Modern Roman}
\usepackage{amsmath,amssymb,mathtools}
\usepackage{unicode-math}
\setmathfont{Latin Modern Math}
\usepackage{enumitem}
\usepackage{xurl,hyperref,bookmark}
\hypersetup{colorlinks=true,urlcolor=blue}
\setcounter{secnumdepth}{0}
\setlength{\parindent}{0pt}
\setlength{\parskip}{5pt}
\setlength{\emergencystretch}{3em}
\begin{document}
\section*{Torsion-Anomalous Conjecture for Semiabelian Varieties}

\subsection{OpenAI's claim}
For subvarieties of abelian varieties over algebraic numbers, the selected manuscript claims finiteness of maximal atypical intersections, with dimensions measured inside the smallest containing torsion coset. This supplies the abelian branch of the semiabelian question. The mathematical statement below specifies that ambient variety and its field of definition.

\subsection{Mathematical statement of the claimed result}
\paragraph{Manuscript 1.}
For a number field $K$, abelian variety $A/K$, and irreducible $X\subseteq A_{\bar K}$, let $S$ be the smallest special subvariety containing $X$; special subvarieties are torsion translates of abelian subvarieties. An irreducible component $Y$ of $X\cap Z$, $Z\subseteq S$ special, is atypical when
\[
 \dim Y>\dim X+\dim Z-\dim S.
\]
The claim is finiteness of maximal atypical subvarieties of $X$, equivalently existence of finitely many proper atypical $Y_j\subsetneq X$ containing every atypical subvariety. No dimension bound on $A$ or $X$ is imposed.

\subsection{Outline of the argument}
For manuscript 1 (family 016; 2026-09-24), a counterexample sequence is decomposed into distinct height scales. Uniform tangent positivity and weighted small sections give a height inequality despite growing homomorphism coefficients and point degrees; passing to suitable quotients rules out the remaining scale configurations.

\subsection{Scope relative to this problem}
An abelian variety is a semiabelian variety with no torus part. The source therefore addresses the abelian case over algebraic numbers, using the smallest containing torsion coset for the atypicality comparison. The displayed question allows an arbitrary semiabelian ambient variety. The manuscript does not assert the extension to a nontrivial torus part or all more general constant fields.

\subsection{Formal verification}
No Lean formalization is listed at this snapshot for the selected manuscript (family 016; 2026-09-24).

This site has not executed the Lean code. Described statement coverage is not an independent proof check.

\subsection{Status caveat}
The manuscripts were produced by an internal OpenAI model. OpenAI cautions that results without formalization can contain mistakes. This site has not independently checked the proofs; the label records the scope of OpenAI's claim.

\subsection{Sources}
\begin{itemize}
\item OpenAI, \emph{The abelian Zilber–Pink conjecture}, 2026-09-24; citation key \texttt{OAI:The-Abelian-Zilber-Pink-Conjecture-September-24-2026}. \href{https://github.com/openai/math/blob/adc7f1241b42e322a6451854ab7e4b4c146bf78a/preprints/The-Abelian-Zilber-Pink-Conjecture-September-24-2026/paper.pdf}{PDF (GitHub)}; \href{https://github.com/openai/math/raw/adc7f1241b42e322a6451854ab7e4b4c146bf78a/preprints/The-Abelian-Zilber-Pink-Conjecture-September-24-2026/paper.pdf}{PDF (direct)}; \href{https://github.com/openai/math/blob/adc7f1241b42e322a6451854ab7e4b4c146bf78a/preprints/The-Abelian-Zilber-Pink-Conjecture-September-24-2026/README.md}{Manuscript page and citation}; \href{https://github.com/openai/math/blob/main/preprints/The-Abelian-Zilber-Pink-Conjecture-September-24-2026/paper.pdf}{Latest version}; Statement source: \href{https://github.com/openai/math/blob/adc7f1241b42e322a6451854ab7e4b4c146bf78a/preprints/The-Abelian-Zilber-Pink-Conjecture-September-24-2026/build/sections/01-introduction.tex\#L43}{thm:main}.
\item \href{https://github.com/openai/math/blob/adc7f1241b42e322a6451854ab7e4b4c146bf78a/CONTENTS.md}{OpenAI catalogue entry}.
\end{itemize}
\end{document}
